\documentclass[11pt]{amsart}

\usepackage[T1]{fontenc}
\usepackage{lmodern}
\usepackage{microtype}
\usepackage{amsmath,amssymb,amsthm}
\usepackage[hidelinks]{hyperref}

\hypersetup{
  pdftitle={The Hafnian van der Waerden Conjecture Fails in Order Six}
}

\newcommand{\haf}{\operatorname{haf}}
\newcommand{\conv}{\operatorname{conv}}

\newtheorem{theorem}{Theorem}
\theoremstyle{remark}
\newtheorem{remark}[theorem]{Remark}

\title[Order-six failure of the hafnian conjecture]
{The Hafnian van der Waerden Conjecture Fails in Order Six}

\date{}

\subjclass[2020]{05C70, 15A15}
\keywords{hafnian, perfect matching polytope, van der Waerden conjecture,
minimal counterexample}

\begin{document}

\begin{abstract}
Friedland formulated a hafnian analogue of the van der Waerden--Tverberg
formula and reported that Gurvits had conjectured its perfect-matching case.
Friedland also records that this case is false once $n$ is sufficiently large,
an observation he credits to S.~Norin. This note locates the least
order at which failure occurs: it is $n=3$, on six vertices. For an explicit
matrix $B\in\Psi_6$,
\[
  \haf(B)=\frac{29}{243}<\frac{3}{25}
  =\haf\!\left(\frac{1}{5}A(K_6)\right),
\]
so the conjectured minimum fails at $n=3$. Failure cannot occur earlier: the
conjecture holds at $n=2$, a case subsumed by Friedland's Lemma~3.1, for which
we include a short direct argument. What is new here is the explicit
order-six matrix and the minimal order it fixes, not the failure of the
conjecture itself. We do not determine the exact minimum on $\Psi_6$.
\end{abstract}

\maketitle

For a perfect matching $M$ of $K_{2n}$, let $A(M)$ denote its symmetric
adjacency matrix, and set
\[
  \Psi_{2n}:=
  \conv\{A(M):M\text{ is a perfect matching of }K_{2n}\}.
\]
This is the set called $\Psi_{2n}$ in \cite{Friedland2011,Friedland2012}, where
it is introduced as the convex set of symmetric doubly stochastic matrices with
zero diagonal whose extreme points are the symmetric permutation matrices with
zero diagonal, and is then described by Edmonds' odd-set inequalities
\cite[Eq.~(1.2)]{Friedland2011}. Those inequalities are part of the
description, so the domain of the conjecture is the perfect matching polytope
displayed above, and not the set cut out by symmetry, nonnegativity, zero
diagonal and double stochasticity alone.
For a symmetric matrix
$C=[c_{ij}]\in\mathbb R_{\ge0}^{2n\times2n}$ with zero diagonal, define
\[
  \haf(C):=
  \sum_M\prod_{\{i,j\}\in M}c_{ij},
\]
where the sum is over the perfect matchings of $K_{2n}$, and let
\[
  \mu_n:=\min_{C\in\Psi_{2n}}\haf(C).
\]
Friedland formulated the corresponding family of hafnian conjectures and
reported that Gurvits had stated the perfect-matching case
\cite[Eq.~(1.4) and the discussion following it]{Friedland2011}; see also
\cite[Eq.~(6.3)]{Friedland2012}. That family is indexed by the size $k$ of the
matchings counted, and the perfect-matching case is $k=n$. In the present
notation, the perfect-matching hafnian van der Waerden conjecture asserts
\begin{equation}\label{eq:conjecture}
  \mu_n=
  \haf\!\left(\frac{1}{2n-1}A(K_{2n})\right),
  \qquad n\ge2.
\end{equation}
That \eqref{eq:conjecture} is false for large $n$ is stated outright in both
printings: ``for $k=n$ and $n$ big enough (1.4) is wrong''
\cite[following Eq.~(1.4)]{Friedland2011}, and likewise for (6.3) in
\cite[\S6]{Friedland2012}. Friedland thanks S.~Norin for pointing out that
failure, and the argument recorded there proceeds through a family of graphs of
Cygan, Pilipczuk and \v{S}krekovski \cite{CPS}; neither that assertion nor its
proof is claimed here.
What follows instead is the least order at which failure occurs, together with
an explicit matrix that witnesses it.

\begin{theorem}\label{thm:main}
Let
\[
B=\frac{1}{9}
\begin{pmatrix}
0&3&3&1&1&1\\
3&0&3&1&1&1\\
3&3&0&1&1&1\\
1&1&1&0&3&3\\
1&1&1&3&0&3\\
1&1&1&3&3&0
\end{pmatrix}.
\]
Then $B\in\Psi_6$ and
\[
  \haf(B)=\frac{29}{243}
  <\frac{3}{25}
  =\haf\!\left(\frac{1}{5}A(K_6)\right).
\]
Consequently, \eqref{eq:conjecture} fails at $n=3$. It holds at $n=2$, a case
subsumed by \cite[Lemma~3.1]{Friedland2011} and reproved directly below, so
$n=3$ is the least admissible value for which it fails.
\end{theorem}

\begin{proof}
Partition $[6]$ into
\[
  T_1=\{1,2,3\},
  \qquad
  T_2=\{4,5,6\}.
\]
For $i\in T_1$ and $j\in T_2$, define the perfect matching
\[
  M_{ij}:=
  \bigl\{\{i,j\},\,T_1\setminus\{i\},\,T_2\setminus\{j\}\bigr\};
\]
the last two sets have cardinality two and hence represent edges. Every edge
between $T_1$ and $T_2$ occurs in exactly one $M_{ij}$, whereas every edge
inside either triangle occurs in exactly three of them. Therefore
\[
  B=\frac{1}{9}
  \sum_{i\in T_1}\sum_{j\in T_2}A(M_{ij}),
\]
so $B\in\Psi_6$.

Every perfect matching of $K_6$ uses an odd number of edges between $T_1$ and
$T_2$, hence either one or three. There are nine matchings with one such edge,
each of weight
\[
  \frac{1}{9}\left(\frac{1}{3}\right)^2=\frac{1}{81},
\]
and six matchings with three such edges, each of weight $1/9^3$. Thus
\[
  \haf(B)
  =9\left(\frac{1}{81}\right)
   +6\left(\frac{1}{729}\right)
  =\frac{29}{243}.
\]
Since $K_6$ has $5!!=15$ perfect matchings,
\[
  \haf\!\left(\frac{1}{5}A(K_6)\right)
  =\frac{15}{5^3}
  =\frac{3}{25},
\]
and
\[
  \frac{3}{25}-\frac{29}{243}
  =\frac{4}{6075}>0.
\]

It remains to verify minimality of the order. Let $M_1,M_2,M_3$ be the three
perfect matchings of $K_4$. Every $C\in\Psi_4$ has the form
\[
  C=\alpha A(M_1)+\beta A(M_2)+\gamma A(M_3),
  \qquad
  \alpha,\beta,\gamma\ge0,
  \qquad
  \alpha+\beta+\gamma=1.
\]
The three matchings partition the edges of $K_4$, and therefore
\[
  \haf(C)=\alpha^2+\beta^2+\gamma^2
  \ge\frac{(\alpha+\beta+\gamma)^2}{3}
  =\frac{1}{3}.
\]
Equality holds precisely at $\alpha=\beta=\gamma=1/3$, which gives
$C=\frac13A(K_4)$ and
\[
  \haf\!\left(\frac13A(K_4)\right)
  =3\left(\frac13\right)^{2}
  =\frac13 ;
\]
so $\mu_2=\frac13$, attained at $\frac13A(K_4)$ and nowhere else. Hence
\eqref{eq:conjecture} holds at $n=2$; this is subsumed by the case $k=2$ of
\cite[Lemma~3.1]{Friedland2011}, and the short direct argument is included only
to keep the minimality half of the theorem self-contained.
\end{proof}

\begin{remark}[Consistency with local minimality]
Friedland proved that $(2n-1)^{-1}A(K_{2n})$ is a strict local minimizer of the
hafnian on $\Psi_{2n}$ \cite[Lemma~3.1]{Friedland2011}.
Theorem~\ref{thm:main} shows that it is already not a global minimizer at
$n=3$. The two are compatible because $B$ is nowhere near
$\frac15A(K_6)$. Indeed a perfect matching of $K_6$ has at most one edge inside
the triple $T_1=\{1,2,3\}$, so every $C=[c_{ij}]\in\Psi_6$ satisfies
$\sum_{\{i,j\}\subseteq T_1}c_{ij}\le1$; this is tight
at $B$, whose entries inside $T_1$ sum to $1$, and strict at $\frac15A(K_6)$,
where they sum to $3/5$. So $B$ lies on the boundary of $\Psi_6$, in a
supporting hyperplane that $\frac15A(K_6)$ misses. Quantitatively, set
$C_t:=(1-t)\frac15A(K_6)+tB$ for $t\in[0,1]$. Each $C_t$ again has constant
entries $(3+2t)/15$ inside the two triples and $(9-4t)/45$ between them, so the
same $9+6$ decomposition gives
\begin{align*}
  \haf(C_t)
  &=\frac{9(9-4t)(3+2t)^2}{45\cdot15^2}
   +\frac{6(9-4t)^3}{45^3}\\
  &=\frac{729+108t^2-112t^3}{6075}
  =\frac{3}{25}+\frac{4t^2(27-28t)}{6075}.
\end{align*}
Thus $\haf(C_t)>\frac3{25}$ for $0<t<\frac{27}{28}$, with equality at
$t=\frac{27}{28}$: the hafnian rises out of $\frac15A(K_6)$, peaks at
$t=\frac9{14}$, and drops below the conjectured value only on the final
$\frac1{28}$ of the segment.
\end{remark}

\begin{remark}[Scope and prior work]\label{rem:scope}
Theorem~\ref{thm:main} refutes \eqref{eq:conjecture} at $n=3$ by an explicit
witness and fixes $n=3$ as the least order at which a witness can exist. It is
not the first refutation of \eqref{eq:conjecture}: both \cite{Friedland2011} and
\cite{Friedland2012} state that the conjecture is wrong for $k=n$ with $n$
large, and no novelty is claimed for that assertion or for the asymptotic
argument behind it. Nor is the case $n=2$ new; it is subsumed by the case $k=2$
of \cite[Lemma~3.1]{Friedland2011}. What is offered here is the matrix $B$, the
exact value $\haf(B)=29/243$, and the resulting identification of $n=3$ as the
minimal order. Theorem~\ref{thm:main} gives only $\mu_3\le29/243$; no equality
claim is made, and the status of the other members of Friedland's family of
hafnian conjectures is not addressed.
\end{remark}

\begin{thebibliography}{3}

\bibitem{CPS}
M.~Cygan, M.~Pilipczuk and R.~\v{S}krekovski,
\emph{A bound on the number of perfect matchings in Klee-graphs},
Discrete Math. Theor. Comput. Sci. \textbf{15} (2013), no.~1, 37--54,
\href{https://doi.org/10.46298/dmtcs.633}{doi:10.46298/dmtcs.633};
cited in \cite{Friedland2011} as a 2009 preprint.

\bibitem{Friedland2011}
S.~Friedland,
\emph{Analogs of the van der Waerden and Tverberg conjectures for haffnians},
\href{https://arxiv.org/abs/1102.2542}{arXiv:1102.2542v2}
[math.CO] (2011).

\bibitem{Friedland2012}
S.~Friedland,
\emph{Results and open problems in matchings in regular graphs},
Electron. J. Linear Algebra \textbf{24} (2012), 18--33,
\href{https://doi.org/10.13001/1081-3810.1577}
{doi:10.13001/1081-3810.1577}.

\end{thebibliography}

\end{document}